\documentclass[11pt,a4paper]{article}

\usepackage{amsmath,amssymb}
\usepackage[utf8]{inputenc}
\usepackage{hyperref}
\usepackage{geometry}
\usepackage{booktabs}
\usepackage{xcolor}
\usepackage{enumitem}

\geometry{margin=1in}

\newcommand{\lean}[1]{\texttt{#1}}
\newcommand{\RH}{\ensuremath{\mathrm{RH}}}

\title{%
  The Fast Typist, Not the Strategist:\\
  What 23 Days of Human-AI Mathematical Collaboration\\
  Reveals About LLMs and Formal Reasoning%
}
\author{
  Jason Robert Gochanour
}
\date{April 19, 2026}

\begin{document}

\maketitle

\begin{abstract}
We report on a 23-day collaboration between a mathematician and an
LLM-based AI coding assistant to formalize the Nyman--Beurling
criterion for the Riemann Hypothesis in Lean~4. The resulting
artifact---the \emph{Cathedral}---comprises 20,000 lines of
machine-verified proof across 78 files, with 644 theorems and 40
axioms. We analyze the division of labor between human and AI,
finding that the AI's primary contribution was \textbf{high-volume
low-creativity work}: Lean syntax generation, Mathlib API discovery,
codebase navigation, and automated refactoring. Mathematical
strategy---proof architecture, axiom classification, and the
discovery of novel techniques---remained exclusively human-driven.
We quantify the AI's productivity gains ($\sim$5$\times$ throughput
vs.\ solo human), identify four failure modes specific to
LLM-assisted theorem proving (stale context, API hallucination,
axiom drift, over-eagerness), and compare with existing AI-for-math
systems (AlphaProof, Lean Copilot, LeanDojo). We argue that current
LLMs are powerful \emph{proof scribes} but not yet
\emph{proof strategists}, and that the bottleneck for AI-assisted
mathematics has shifted from code generation to mathematical taste.
\end{abstract}

\tableofcontents

% ================================================================
\section{Introduction}
% ================================================================

Can AI prove theorems? The question has animated research since
Newell and Simon's Logic Theorist (1956) and has recently gained
momentum with AlphaProof's IMO silver medal (2024), Lean Copilot
(2024), and LeanDojo's retrieval-augmented proving (2023).

These systems operate at the level of individual lemmas: given a
statement, find its proof. The Cathedral project operates at a
qualitatively different scale: not proving one theorem, but designing
an entire proof \emph{architecture}---78 files, 11 modules, 644
theorems---for a major open problem in mathematics.

This paper analyzes the human-AI collaboration that produced the
Cathedral, focusing on the question: \emph{what did the AI actually
do?}

\subsection{Summary of Findings}

\begin{enumerate}
  \item The AI generated $\sim$85\% of all Lean code by volume, but
    made $\sim$0\% of architectural decisions.
  \item The AI's primary value was in \textbf{throughput}: what would
    take a human mathematician hours of Lean syntax wrestling took
    the AI minutes.
  \item The AI failed consistently at tasks requiring
    \textbf{mathematical taste}: choosing between proof strategies,
    identifying which axioms to eliminate first, and recognizing
    when an approach was aesthetically wrong even if logically valid.
  \item The collaboration produced a 54\% discard rate, higher than
    either human or AI would produce alone, suggesting that the
    pairing enables more aggressive exploration.
\end{enumerate}

% ================================================================
\section{System and Methodology}
\label{sec:system}
% ================================================================

\subsection{The AI System}

The AI assistant was an LLM-based coding agent with the following
capabilities:

\begin{itemize}
  \item File reading and writing within the project workspace.
  \item Terminal command execution (\lean{lake build}, \lean{lake env lean}).
  \item Web search for Mathlib documentation and API references.
  \item Context window of $\sim$200K tokens (insufficient for the
    full 20,000-line codebase).
  \item No persistent memory between sessions (context rebuilt each
    time).
\end{itemize}

The AI had no special-purpose theorem-proving capabilities. It was
a general-purpose coding assistant that happened to be writing Lean~4.

\subsection{The Human}

The human collaborator is a mathematician with expertise in analytic
number theory and experience with formal verification. The human
provided:

\begin{itemize}
  \item Mathematical strategy and proof architecture.
  \item Axiom classification and elimination priorities.
  \item Quality control (deciding when a proof was ``done'' vs.\
    ``correct but ugly'').
  \item Error diagnosis when the AI's approach failed.
\end{itemize}

\subsection{Session Protocol}

Development occurred in sessions of 2--8 hours over 23 days. A
typical session:

\begin{enumerate}
  \item \textbf{Human sets goal}: ``Eliminate axiom X'' or ``Prove
    theorem Y.''
  \item \textbf{AI reads context}: Relevant source files, dependencies,
    Mathlib API.
  \item \textbf{AI proposes strategy}: ``I'll use Mathlib's
    \lean{riemannZeta\_ne\_zero\_of\_one\_le\_re} and the functional
    equation.''
  \item \textbf{Human approves or redirects}: ``Good approach'' or
    ``No, try the contrapositive via the Mellin transform instead.''
  \item \textbf{AI writes code}: Full Lean proof, typically 50--300
    lines.
  \item \textbf{Compiler checks}: AI runs \lean{lake env lean} and
    iterates on errors.
  \item \textbf{Regression}: AI runs \lean{lake build} to check
    full-codebase consistency.
  \item \textbf{Human reviews}: Checks mathematical content,
    docstrings, axiom impact.
\end{enumerate}

% ================================================================
\section{Division of Labor: Quantitative Analysis}
\label{sec:division}
% ================================================================

\subsection{Task Taxonomy}

We classify all development tasks into five categories:

\begin{center}
\begin{tabular}{@{}llcc@{}}
\toprule
\textbf{Category} & \textbf{Example} & \textbf{Human} & \textbf{AI} \\
\midrule
Strategy & ``Use Rank-1 Mellin for converse'' & 100\% & 0\% \\
Architecture & ``Create 11-module structure'' & 90\% & 10\% \\
Implementation & ``Write the Lean proof'' & 15\% & 85\% \\
Maintenance & ``Update docstrings, fix imports'' & 5\% & 95\% \\
Debugging & ``Why does this type not unify?'' & 40\% & 60\% \\
\bottomrule
\end{tabular}
\end{center}

\subsection{Volume Analysis}

\begin{center}
\begin{tabular}{@{}lr@{}}
\toprule
\textbf{Metric} & \textbf{Value} \\
\midrule
Total LOC produced (active + archived) & 43,041 \\
Estimated AI-authored LOC & $\sim$36,500 (85\%) \\
Estimated human-authored LOC & $\sim$6,500 (15\%) \\
Files created & 174 \\
Files archived (discarded) & 96 (55\%) \\
Axioms declared & $\sim$80 (over lifetime) \\
Axioms eliminated & 40 \\
Novel mathematical insights & 5 \\
Of which AI-originated & 0 \\
\bottomrule
\end{tabular}
\end{center}

The most striking number: \textbf{zero} novel mathematical insights
originated from the AI. Every creative idea---the Rank-1 Mellin
Miracle, the Parseval Bridge, the axiom tiering system, the Bartlett
window interpretation---came from the human. The AI implemented them.

\subsection{Throughput Analysis}

Estimated time comparison for key tasks:

\begin{center}
\begin{tabular}{@{}lrrr@{}}
\toprule
\textbf{Task} & \textbf{Human solo} & \textbf{With AI} & \textbf{Speedup} \\
\midrule
Write 200-line Lean proof & 4--8 hours & 30--60 min & $\sim$8$\times$ \\
Find correct Mathlib API & 1--2 hours & 5--15 min & $\sim$6$\times$ \\
Debug type unification error & 30--60 min & 10--20 min & $\sim$3$\times$ \\
Update all docstrings & 2--3 hours & 15--30 min & $\sim$8$\times$ \\
Design proof architecture & 2--4 hours & 2--4 hours & $1\times$ \\
Choose next axiom to eliminate & 30 min & 30 min & $1\times$ \\
\bottomrule
\end{tabular}
\end{center}

The pattern is clear: the AI accelerates \textbf{execution} but not
\textbf{decision-making}. For ``what to do'' questions, the AI
provides no speedup. For ``how to do it'' questions, the speedup is
3--8$\times$.

% ================================================================
\section{AI Failure Modes}
\label{sec:failures}
% ================================================================

We identify four recurring failure modes:

\subsection{Stale Context}

The AI's context window ($\sim$200K tokens) cannot hold the full
codebase ($\sim$20K lines of active code + Mathlib). At the start
of each session, the AI must re-read relevant files, occasionally
missing dependencies or using stale mental models.

\textbf{Symptom}: The AI writes code that references a function
signature that was changed in a previous session.

\textbf{Frequency}: $\sim$1--2 times per session.

\textbf{Mitigation}: The AI was trained to always re-read touched
files before modifying them.

\subsection{API Hallucination}

The AI generates plausible-sounding but non-existent Mathlib function
names. For example, suggesting \lean{Integrable.sum\_le\_integral\_Ico}
when the actual name is \lean{AntitoneOn.inner\_le\_lnorm\_mul\_lnorm}.

\textbf{Symptom}: Compilation failure with ``unknown identifier.''

\textbf{Frequency}: $\sim$3--5 times per proof attempt.

\textbf{Mitigation}: The AI was trained to grep the Mathlib source
for function names before using them, and to propose ``search
strategies'' rather than specific API calls.

\subsection{Axiom Drift}

Adding a new file or import can silently pull axioms into the
transitive dependency graph. The AI, not tracking axiom counts
proactively, sometimes introduced unnecessary axioms.

\textbf{Symptom}: \lean{\#print axioms} shows unexpected axioms on
the crown path.

\textbf{Frequency}: $\sim$1 time per major refactoring.

\textbf{Mitigation}: The human established a policy of running
\lean{\#print axioms} after every structural change.

\subsection{Over-Eagerness}

The AI tends to produce a ``working'' proof quickly but without
considering elegance, minimality, or future maintainability. A proof
that uses 5 axioms when 2 would suffice is ``correct'' but not
``good.''

\textbf{Symptom}: The human reviews a proof and says ``This is
right but wrong.''

\textbf{Frequency}: Persistent, requiring constant human oversight.

\textbf{Mitigation}: The human established a ``sweep the floor''
maintenance policy and reviewed all proofs for axiom minimality.

% ================================================================
\section{Comparison with Existing AI-for-Math Systems}
\label{sec:comparison}
% ================================================================

\begin{center}
\small
\begin{tabular}{@{}lcccc@{}}
\toprule
& \textbf{AlphaProof} & \textbf{LeanDojo} & \textbf{Lean Copilot} & \textbf{Cathedral} \\
\midrule
Scale & Single problem & Single lemma & Single tactic & 78-file project \\
Autonomy & High & High & Low & Low (human-directed) \\
Math level & IMO & Mathlib & Mathlib & Research frontier \\
Novel math & Some & None & None & None (from AI) \\
Architecture & N/A & N/A & N/A & Human-designed \\
LOC & $\sim$100s & $\sim$10s & $\sim$1s & 20,000 \\
\bottomrule
\end{tabular}
\end{center}

The Cathedral occupies a unique niche: large-scale, human-directed,
operating at the research frontier. Existing AI-for-math systems
excel at small-scale autonomous proving (AlphaProof) or tactic
suggestion (Lean Copilot). None have been applied to a multi-file
formal verification project at the scale of the Cathedral.

\subsection{What Would Change the Game}

For AI to move from ``fast typist'' to ``proof strategist,'' the
following capabilities would be needed:

\begin{enumerate}
  \item \textbf{Long-term memory}: The ability to maintain a coherent
    model of a 20,000-line codebase across sessions.
  \item \textbf{Axiom awareness}: Built-in tracking of which axioms
    a proof depends on, with optimization for minimality.
  \item \textbf{Mathematical taste}: The ability to distinguish
    between ``correct but ugly'' and ``correct and elegant'' proofs.
    This is the hardest capability to define or train for.
  \item \textbf{Strategic planning}: The ability to decompose a
    multi-week project into phases, choose proof architectures, and
    backtrack when an approach fails.
\end{enumerate}

Of these, (1) is a systems problem, (2) is a tools problem, (3) is
an open research question, and (4) is the frontier.

% ================================================================
\section{The 54\% Discard Rate}
\label{sec:discard}
% ================================================================

The Cathedral produced 174 files, of which 96 (55\%) were eventually
archived. This discard rate is unusually high for software engineering
but may be characteristic of formal proof search.

\subsection{Why So Much Waste?}

\begin{itemize}
  \item \textbf{Exploration is cheap with AI}: Because the AI can
    produce a 200-line proof attempt in 30 minutes, there is little
    cost to trying an approach that might fail. Without the AI, each
    attempt would take 4--8 hours, strongly discouraging exploration.
  \item \textbf{Proof search is inherently exploratory}: Unlike
    software, where the target behavior is known, proof search
    explores an unknown space. Dead ends are information.
  \item \textbf{The archive preserves negative knowledge}: Each
    archived file documents an approach that doesn't work, preventing
    future explorers from repeating it.
\end{itemize}

\subsection{Is the Discard Rate Optimal?}

We conjecture that the optimal discard rate for a frontier proof
project lies between 30\% and 70\%. Below 30\%, the exploration is
too conservative; above 70\%, the implementation lacks focus. The
Cathedral's 55\% suggests a roughly balanced exploration-exploitation
tradeoff, enabled by the AI's ability to make exploration cheap.

% ================================================================
\section{Lessons for AI-Assisted Mathematics}
\label{sec:lessons}
% ================================================================

\begin{enumerate}
  \item \textbf{AI multiplies throughput, not insight.} The
    Cathedral could not have been built in 23 days without AI
    assistance. It also could not have been built \emph{at all}
    without human mathematical insight. The AI was a
    $5\times$ throughput multiplier, not a creativity enhancer.

  \item \textbf{The bottleneck is taste, not skill.} The AI can
    write correct Lean code faster than any human. What it cannot
    do is decide \emph{which} correct code to write. Mathematical
    taste---the ability to distinguish elegant from merely valid
    approaches---is the binding constraint.

  \item \textbf{Verification is the trust anchor.} In a workflow
    where AI generates most of the code, the Lean type-checker
    serves as the ultimate quality gate. The human doesn't need to
    verify every line---the compiler does that. The human needs to
    verify that the \emph{right thing} was proved.

  \item \textbf{Discard culture is healthy.} The willingness to
    archive 96 files without regret is essential. Sunk cost fallacy
    is the enemy of proof search. The AI's ability to regenerate
    code cheaply makes discarding psychologically easier.

  \item \textbf{The AI is best as a scribe, not a sage.} The most
    productive interaction pattern was: human provides mathematical
    direction in natural language, AI translates to Lean. The least
    productive was: human asks AI for mathematical ideas. The AI's
    mathematical suggestions were generic, often wrong, and never
    novel.
\end{enumerate}

% ================================================================
\section{Conclusion}
% ================================================================

The Cathedral demonstrates that LLMs can meaningfully accelerate
frontier mathematics when used as proof scribes under human direction.
The collaboration produced 20,000 lines of machine-verified proof in
23 days---a pace that would be extraordinary for a solo mathematician
working with Lean~4.

But the AI contributed no new mathematical ideas. Every creative
insight---the Rank-1 Mellin Miracle, the Parseval Bridge, the
axiom tiering system---came from the human. The AI's contribution
was translation, not creation; velocity, not direction.

The title of this paper is our conclusion: the AI is the fast
typist, not the strategist. This is not a criticism. Fast typing is
extraordinarily valuable when the strategist knows where to go.
The next cathedral will be built faster. But it will still need a
human to draw the blueprints.

\vspace{0.5cm}
\noindent\textbf{Source code:} \url{https://github.com/jrgochan/prime}

\begin{thebibliography}{99}

\bibitem{alphaproof}
Google DeepMind, \emph{AlphaProof and AlphaGeometry 2}, 2024.

\bibitem{leandojo}
K.~Yang et al., \emph{LeanDojo: Theorem Proving with Retrieval-Augmented
Language Models and the LeanDojo Benchmark}, NeurIPS 2023.

\bibitem{leancopilot}
S.~Song et al., \emph{Towards Large Language Models as Copilots for
Theorem Proving in Lean}, 2024.

\bibitem{newell1956}
A.~Newell and H.~Simon, \emph{The Logic Theory Machine},
IRE Transactions on Information Theory, 1956.

\bibitem{lean4}
L.~de~Moura et al., \emph{The Lean~4 Theorem Prover}, CADE-28, 2021.

\end{thebibliography}

\end{document}
